\documentclass{article}
\usepackage{graphicx} % Required for inserting images
\usepackage{amsmath, amssymb, amsthm, graphicx, geometry}
\geometry{margin=1in}

\maketitle
\author{Pablo Nogueira Grossi}
\date{March 2026}

\maketitle\section{Scale-Invariant Development Levels in the \texorpdfstring{$g$}{g}-Series}

In this section we formalise the notion of development levels in the
$g$–series operator algebra, without committing to any specific
numerical index such as $33$. The goal is to separate (i) the
\emph{mathematical structure} of development levels from (ii) any
\emph{empirical} observation about which index real systems tend to
occupy.

\subsection{State space, scales, and the \texorpdfstring{$g$}{g}-series}

\begin{definition}[dm3 state space]
A \emph{dm3 system} is a triple


\[
D = (M, \Gamma, V)
\]


where $M$ is a smooth Riemannian $3$–manifold, $\Gamma \subset M$ is a
hyperbolic limit cycle, and $V : M \to \mathbb{R}$ is a $C^2$ Lyapunov
function satisfying the eight dm3 axioms of Generative Contact
Mechanics (existence of $\Gamma$, orbital stability, stochastic
stability, etc.).
\end{definition}

\begin{definition}[Scale]
A \emph{scale} is an equivalence class of dm3 systems under smooth
rescalings of space, time, and Lyapunov potential that preserve:
\begin{enumerate}
  \item the existence and hyperbolicity of $\Gamma$,
  \item the sign of the Lyapunov exponent $\mu_{\max}$,
  \item the sign of the noise threshold $\tau$.
\end{enumerate}
We write $[D]$ for the scale containing a dm3 system $D$.
\end{definition}

\begin{definition}[$g$–series at fixed scale]
Fix a scale $[D]$ and choose a representative $D_0 \in [D]$. A
\emph{$g$–series} at this scale is a sequence


\[
\{D_n\}_{n \in \mathbb{N}}, \qquad D_n \in [D],
\]


together with operators


\[
g_n : [D] \to [D], \qquad n \in \mathbb{N},
\]


such that $D_{n+1} = g_n(D_n)$ and each $g_n$ is built from the
GCM operator grammar (compositions of $g$–, $L$–, $R$–, and $U$–operators)
and preserves the dm3 axioms.
\end{definition}

\subsection{Full development and dynamic center}

We now define two structural notions: full development at a given
scale, and the dynamic center of that scale. Neither involves any
specific numerical index.

\begin{definition}[Development functional]
Let $[D]$ be a fixed scale. A \emph{development functional} at this
scale is a map


\[
\mathcal{D} : [D] \to [0,1]
\]


such that:
\begin{enumerate}
  \item $\mathcal{D}(D) = 0$ if and only if $D$ is minimally developed
  (e.g.\ smallest admissible basin, smallest admissible $\tau$);
  \item $\mathcal{D}(D) = 1$ if and only if $D$ is maximally developed
  at this scale (no further admissible operator in the grammar can
  increase $\tau$ or the basin without crossing a boundary $B_2$ or
  $B_3$);
  \item $\mathcal{D}$ is non-decreasing under admissible $g$–moves:
  if $D' = g(D)$ for some grammar operator $g$, then
  $\mathcal{D}(D') \geq \mathcal{D}(D)$.
\end{enumerate}
\end{definition}

\begin{definition}[Full development level]
Let $\{D_n\}_{n \in \mathbb{N}}$ be a $g$–series at scale $[D]$ and
$\mathcal{D}$ a development functional. A \emph{full development level}
is an index $n_{\mathrm{fd}} \in \mathbb{N}$ such that:
\begin{enumerate}
  \item $\mathcal{D}(D_{n_{\mathrm{fd}}}) = 1$;
  \item for all $m \geq n_{\mathrm{fd}}$ with $D_m$ still in the same
  scale $[D]$, we have $\mathcal{D}(D_m) = 1$.
\end{enumerate}
We then say that $D_{n_{\mathrm{fd}}}$ is \emph{fully developed at
scale $[D]$}.
\end{definition}

\begin{definition}[Dynamic center]
Let $\{D_n\}_{n \in \mathbb{N}}$ and $\mathcal{D}$ be as above. A
\emph{dynamic center} index $n_{\mathrm{c}}$ is an index such that:
\begin{enumerate}
  \item $\mathcal{D}(D_{n_{\mathrm{c}}}) \in (0,1)$;
  \item for all sufficiently small perturbations of the environment
  (noise level, coupling strength) that keep the system within the
  same scale $[D]$, the induced evolution of the $g$–series returns
  to a neighborhood of $D_{n_{\mathrm{c}}}$;
  \item $n_{\mathrm{c}}$ is not a fixed point of the dynamical flow in
  state space, but is a fixed point of the \emph{grammar}: the local
  operator choices that restore coherence and prevent collapse
  (the $L$– and $R$–operators) tend to steer the system back toward
  $D_{n_{\mathrm{c}}}$.
\end{enumerate}
\end{definition}

\begin{remark}
In applications, one often normalises $\mathcal{D}$ so that
$\mathcal{D}(D_{n_{\mathrm{c}}}) = \tfrac{1}{2}$, i.e.\ the dynamic
center sits halfway between minimal and full development at the
current scale. This is a choice of scale for $\mathcal{D}$, not a
mathematical necessity.
\end{remark}

\subsection{Scale crossing and seed states}

\begin{definition}[Scale crossing]
Let $[D]$ and $[D']$ be two scales with $[D']$ strictly larger (in the
sense of additional degrees of freedom or larger ambient manifold).
A \emph{scale crossing} from $[D]$ to $[D']$ occurs when an admissible
operator in the grammar pushes a fully developed state
$D_{n_{\mathrm{fd}}} \in [D]$ across a boundary $B_2$ or $B_3$,
producing a new dm3 system $D'_0 \in [D']$ that satisfies the eight
axioms at the higher scale.
\end{definition}

\begin{definition}[Seed state at higher scale]
The system $D'_0 \in [D']$ obtained by scale crossing from a fully
developed $D_{n_{\mathrm{fd}}} \in [D]$ is called a \emph{seed state}
for scale $[D']$. It is minimally developed at $[D']$ in the sense that
$\mathcal{D}'(D'_0)$ is close to $0$ for any development functional
$\mathcal{D}'$ at scale $[D']$.
\end{definition}

\begin{proposition}[Scale-invariant development structure]
Let $[D]$ be a scale, $\{D_n\}_{n \in \mathbb{N}}$ a $g$–series at this
scale, and $\mathcal{D}$ a development functional. Assume:
\begin{enumerate}
  \item There exists a full development level $n_{\mathrm{fd}}$.
  \item There exists a dynamic center index $n_{\mathrm{c}}$.
  \item A scale crossing from $[D]$ to some $[D']$ occurs at or after
  $n_{\mathrm{fd}}$.
\end{enumerate}
Then:
\begin{enumerate}
  \item The image of $D_{n_{\mathrm{fd}}}$ under scale crossing is a
  seed state $D'_0 \in [D']$.
  \item There exists a $g$–series $\{D'_m\}_{m \in \mathbb{N}}$ at
  scale $[D']$ and a development functional $\mathcal{D}'$ such that
  the same qualitative structure holds: a full development level
  $n'_{\mathrm{fd}}$ and a dynamic center $n'_{\mathrm{c}}$.
\end{enumerate}
In particular, the notions of full development and dynamic center are
\emph{scale-invariant}: they recur at every scale where the dm3 axioms
and the operator grammar are valid.
\end{proposition}

\begin{proof}
The existence of $D'_0$ as a seed state follows from the definition of
scale crossing and the dm3 axioms at the higher scale. The construction
of $\{D'_m\}$ and $\mathcal{D}'$ repeats the same operator grammar and
development functional construction at $[D']$, using the fact that the
axioms are scale-free. The qualitative properties of full development
and dynamic center are preserved by construction. 
\end{proof}

\subsection{Empirical indices and the number \texorpdfstring{$33$}{33}}

\begin{remark}[Empirical development index]
In numerical experiments, one may observe that for a particular
discretisation, operator pipeline, and persistent-homology readout,
the dynamic center index $n_{\mathrm{c}}$ takes a specific value (for
example $n_{\mathrm{c}} = 33$). This is an \emph{empirical constant}
of that experimental setup, not a mathematical invariant of the
theory. The formalism above deliberately avoids baking any such
numerical value into the definitions.
\end{remark}

\begin{conjecture}[Example of an empirical pattern]
For a specific class of discretised dm3 systems and a fixed
persistent-homology pipeline, the dynamic center index
$n_{\mathrm{c}}$ is observed to be $33$ and is robust under a family
of perturbations. Explaining this robustness from first principles
is an open problem.
\end{conjecture}
